\section{Support-relation properties and controlled validation}
\label{sec:s-support}

\subsection{Properties}

The support relation treats benchmark reporting as a measurement problem: claims, operational interpretations, and evidence scopes are recorded together~\cite{jacobs2021measurement,reuel2024betterbench,bordes2025evalfacts}. Table~\ref{tab:s-support-properties} summarizes the five properties used by the validator.

\begin{table}[tbp]
\centering
\footnotesize
\setlength{\tabcolsep}{3.2pt}
\renewcommand{\arraystretch}{1.12}
\caption{Properties enforced by the typed support relation.}
\label{tab:s-support-properties}
\begin{tabularx}{\textwidth}{@{}l >{\raggedright\arraybackslash}p{0.25\textwidth} Y@{}}
\toprule
\textbf{ID} & \textbf{Property} & \textbf{Operational consequence} \\
\midrule
P1 & Scope monotonicity & Widening a claim adds requirements; it cannot erase requirements already named. \\
P2 & Evidence restriction & Removing the only admissible witness for a required element breaks support. \\
P3 & Rank invariance & Strictly increasing score calibration preserves AUROC/AUPRC ordering but may change a fixed-threshold decision. \\
P4 & Protocol non-substitutability & A result under one claim-relevant visibility or selection coordinate cannot fill another without an invariance argument. \\
P5 & Resource separation & A resource-blocked cell is measurable for feasibility but unordered by predictive performance. \\
\bottomrule
\end{tabularx}
\vspace{2pt}
\begin{minipage}{\textwidth}\footnotesize\emph{Scope and reading.} These are specification and ordering properties, not claims that the ontology is complete or that the validator establishes scientific truth.\end{minipage}
\end{table}


\begin{proposition}[Scope monotonicity of requirements]
If \(c'\) widens \(c\) by adding datasets, protocols, models, seeds, or resource regimes while holding the predicate fixed, then \(\Required(c)\subseteq\Required(c')\).
\end{proposition}
\begin{proof}[Proof sketch]
The requirement builder retains every named cell and auxiliary record and adds the widened scope. It never deletes an existing requirement. The property concerns requirements; it does not make truth monotone.
\end{proof}

\begin{proposition}[Evidence restriction]
If \(e\) is the only admissible witness for an element of \(\Required(c)\), then removing \(e\) breaks support for \(c\).
\end{proposition}
\begin{proof}
The required element becomes uncovered. Duplicate files with the same validated scientific key do not create a second independent witness.
\end{proof}

\begin{proposition}[Rank invariance under increasing calibration]
If \(g\) is strictly increasing and tie order is unchanged, replacing \(q_i\) with \(g(q_i)\) preserves AUROC and AUPRC but may change decisions at a fixed numerical threshold.
\end{proposition}
\begin{proof}
The rank metrics depend on score order, which \(g\) preserves. The selected set \(\{i:q_i\geq\tau\}\) need not equal \(\{i:g(q_i)\geq\tau\}\).
\end{proof}

\begin{proposition}[Protocol non-substitutability]
A result under \(\Pi_1\) cannot satisfy a requirement scoped to \(\Pi_2\) when the contracts differ on a claim-relevant coordinate, absent a separate invariance argument.
\end{proposition}
\begin{proof}[Proof sketch]
The contract is part of the scientific key. Changing a claim-relevant coordinate changes admissible information or the evaluation condition.
\end{proof}

\begin{proposition}[Resource separation]
A resource-blocked cell has no predictive ordering relation to a measured cell.
\end{proposition}
\begin{proof}
The blocked cell has no observed AUPRC, AUROC, or F1. It can be compared only as a feasibility outcome under \(\mathcal{R}\).
\end{proof}

\subsection{Controlled implementation validation}

Table~\ref{tab:s-validation-cases} reports the complete controlled suite. Each case starts from a frozen claim/evidence relation and changes one required property in a sandbox: scope, seed completeness, prediction presence, integrity, construct identity, resource state, or direction.

\begingroup
\footnotesize
\setlength{\tabcolsep}{3pt}
\renewcommand{\arraystretch}{1.10}
\begin{longtable}{@{}p{0.07\textwidth} p{0.18\textwidth} p{0.38\textwidth} p{0.15\textwidth} p{0.15\textwidth}@{}}
\caption{Controlled support-relation validation cases (14 of 14 matched the expected status).}\label{tab:s-validation-cases}\\
\toprule
\textbf{Case} & \textbf{Base claim} & \textbf{Controlled mutation} & \textbf{Expected} & \textbf{Observed} \\
\midrule
\endfirsthead
\caption[]{Controlled support-relation validation cases (14 of 14 matched the expected status). (continued)}\\
\toprule
\textbf{Case} & \textbf{Base claim} & \textbf{Controlled mutation} & \textbf{Expected} & \textbf{Observed} \\
\midrule
\endhead
\midrule
\multicolumn{5}{r}{\footnotesize Continued on next page}\\
\endfoot
\bottomrule
\endlastfoot
FV01 & C01 Elliptic visibility ranking & none; complete locked grid & \texttt{SUPPORTED} & \texttt{SUPPORTED} \\
FV02 & C01 Elliptic visibility ranking & remove one model/protocol cell in sandbox & \texttt{BLOCKED\_\allowbreak{}INCOMPLETE\_\allowbreak{}SCOPE} & \texttt{BLOCKED\_\allowbreak{}INCOMPLETE\_\allowbreak{}SCOPE} \\
FV03 & C01 Elliptic visibility ranking & remove seed 10 from one cell in sandbox & \texttt{BLOCKED\_\allowbreak{}INCOMPLETE\_\allowbreak{}SEEDS} & \texttt{BLOCKED\_\allowbreak{}INCOMPLETE\_\allowbreak{}SEEDS} \\
FV04 & C01 Elliptic visibility ranking & erase one prediction-manifest reference in sandbox & \texttt{BLOCKED\_\allowbreak{}MISSING\_\allowbreak{}PREDICTIONS} & \texttt{BLOCKED\_\allowbreak{}MISSING\_\allowbreak{}PREDICTIONS} \\
FV05 & C10 Small GINE AUPRC result & widen scope from Small to Medium & \texttt{RESOURCE\_\allowbreak{}BLOCKED} & \texttt{RESOURCE\_\allowbreak{}BLOCKED} \\
FV06 & C05 IBM AML baseline grid & widen scope from Small/Medium to all official variants & \texttt{RESOURCE\_\allowbreak{}BLOCKED} & \texttt{RESOURCE\_\allowbreak{}BLOCKED} \\
FV07 & V22 loss robustness & substitute a partial failed import for a full10 lane & \texttt{EXCLUDED\_\allowbreak{}INTEGRITY} & \texttt{EXCLUDED\_\allowbreak{}INTEGRITY} \\
FV08 & V24 temporal stress & treat three metadata labels as distinct temporal windows & \texttt{EXCLUDED\_\allowbreak{}CONSTRUCT\_\allowbreak{}INVALID} & \texttt{EXCLUDED\_\allowbreak{}CONSTRUCT\_\allowbreak{}INVALID} \\
FV09 & DGraphFin fixed GAT performance & promote the T4-OOM cell to a predictive comparison & \texttt{RESOURCE\_\allowbreak{}BLOCKED} & \texttt{RESOURCE\_\allowbreak{}BLOCKED} \\
FV10 & DGraphFin GraphSAGE max-pool rerun & promote runbook/status files without imported rows & \texttt{RESOURCE\_\allowbreak{}BLOCKED} & \texttt{RESOURCE\_\allowbreak{}BLOCKED} \\
FV11 & IBM sender-receiver construction & count the contract alias as an independent replication & \texttt{EXCLUDED\_\allowbreak{}CONSTRUCT\_\allowbreak{}INVALID} & \texttt{EXCLUDED\_\allowbreak{}CONSTRUCT\_\allowbreak{}INVALID} \\
FV12 & C04 universal graph harm & widen mixed scoped results into a universal directional claim & \texttt{REFUTED\_\allowbreak{}IN\_\allowbreak{}SCOPE} & \texttt{REFUTED\_\allowbreak{}IN\_\allowbreak{}SCOPE} \\
FV13 & C17 universal GraphSafe improvement & widen bounded mixed case study to universal dominance & \texttt{REFUTED\_\allowbreak{}IN\_\allowbreak{}SCOPE} & \texttt{REFUTED\_\allowbreak{}IN\_\allowbreak{}SCOPE} \\
FV14 & IBM early-to-late visibility & describe the shared first-60\% label-free node-history map as first-50\% training-only features & \texttt{EXCLUDED\_\allowbreak{}CONSTRUCT\_\allowbreak{}INVALID} & \texttt{EXCLUDED\_\allowbreak{}CONSTRUCT\_\allowbreak{}INVALID} \\
\end{longtable}
\noindent\footnotesize\emph{Scope and reading.} The cases localize incomplete scope, missing seeds or predictions, integrity and construct failures, resource blocks, and directional refutation. They are project-authored conformance tests, not external validation.\par
\endgroup


All 14 observed statuses equal their expected statuses. FV08 rejects the V24 stress-label interpretation because the label never reached the harness; FV11 rejects the sender-receiver alias as an independent method; FV14 rejects a first-50\%-only description of IBM account histories. These cases show that the implementation preserves the declared status distinctions.

\begin{figure}[tbp]
    \centering
    \includegraphics[width=0.86\textwidth]{fig08_claim_support_validation.pdf}
    \caption{Observed transitions for the 14 controlled claim/evidence mutations. The suite tests conformance to the project specification, not universal validator correctness.}
    \label{fig:s-support-validation-v2}
\end{figure}

\paragraph{What this does not establish.}
The schema, cases, and implementation were authored in the same project. Passing the suite does not prove ontology completeness, scientific truth, causal validity, curator agreement, fairness, or responsible deployment.
